\documentclass[11pt]{article}

\usepackage[margin=1in]{geometry}
\usepackage{amsmath, amssymb}
\usepackage{booktabs}
\usepackage{array}
\usepackage{xcolor}
\usepackage{hyperref}
\usepackage{microtype}
\usepackage[strings]{underscore}

\newcommand{\code}[1]{\texttt{#1}}
\newcommand{\R}{\mathbb{R}}

\hypersetup{
    colorlinks=true,
    linkcolor=blue!50!black,
    urlcolor=blue!50!black,
    citecolor=blue!50!black,
    pdftitle={Hessian Spectral Density Estimation: Formalism, Implementation, Observables},
}

\title{Hessian Spectral Density Estimation:\\
Formalism, Implementation, and Observables}
\author{eigengrokking}
\date{\today}

\begin{document}
\maketitle

\begin{abstract}
This note documents the spectral analysis module (\code{analysis/eigenthings.py},
\code{analysis/spectral_observables.py}) used to track the loss Hessian's
eigenvalue distribution over training. It covers, in order: the mathematical
formalism (Hessian-vector products, Lanczos tridiagonalization, Stochastic
Lanczos Quadrature), how the code implements that formalism, the eight scalar
observables extracted per checkpoint and their exact code-level definitions,
and how the two governing hyperparameters --- Lanczos depth \(m\) and probe
count \(k\) --- trade off against estimation accuracy for each observable.
The accuracy claims in Section~\ref{sec:accuracy} are empirical, taken from
\code{reports/spectral_validation.md}, which validates the estimator against
two exact-reference fixtures; see that document for the underlying sweeps.
\end{abstract}

\tableofcontents

% ============================================================
\section{Mathematical formalism}
\label{sec:formalism}

\subsection{Setup: the Hessian and its spectral measure}

Let \(\theta \in \R^n\) be the trainable parameters of the model (flattened
across all tensors) and \(\mathcal{L}(\theta)\) the scalar loss evaluated on
a fixed batch. The object of interest is the Hessian
\[
H \;=\; \nabla^2_\theta \mathcal{L}(\theta) \;\in\; \R^{n\times n},
\]
a real symmetric matrix with eigendecomposition \(H = V\Lambda V^\top\),
eigenvalues \(\lambda_1 \le \cdots \le \lambda_n\). Since \(H\) is
\emph{not} a Gauss--Newton or Fisher approximation, it is generically
indefinite: negative eigenvalues correspond to directions of negative
curvature at \(\theta\).

For \(n\) in the \(10^3\)--\(10^7\) range, forming \(H\) explicitly is
infeasible. The quantity we actually want is the \emph{empirical spectral
measure}
\[
\mu_H \;=\; \frac{1}{n}\sum_{i=1}^n \delta(\lambda - \lambda_i),
\]
which we only ever access indirectly, through two primitives: (i)
Hessian-vector products \(v \mapsto Hv\), and (ii) an algorithm --
Stochastic Lanczos Quadrature (SLQ) -- that turns a modest number of HVPs
into an estimate of \(\mu_H\) and of linear functionals of it (trace,
tail mass, entropy, \dots).

\subsection{Hessian-vector products via double backpropagation}
\label{sec:hvp}

Pearlmutter's trick computes \(Hv\) exactly, without materializing \(H\),
by differentiating the loss twice. Writing \(g(\theta) = \nabla_\theta
\mathcal{L}(\theta)\),
\[
Hv \;=\; \nabla_\theta\!\left(g(\theta)^\top v\right),
\]
because \(\nabla_\theta (g^\top v) = \nabla_\theta\!\left(v^\top \nabla_\theta
\mathcal{L}\right) = (\nabla^2_\theta \mathcal{L}) v = Hv\) (\(v\) held
fixed, not a function of \(\theta\)). Concretely: one backward pass with
\code{create_graph=True} produces \(g\) as a differentiable tensor; a second
backward pass through the scalar \(g^\top v\) yields \(Hv\). This is the
\emph{true} Hessian, not a Gauss--Newton/Fisher surrogate ---
important here since the surrogate is positive semi-definite by
construction and would zero out exactly the negative-curvature signal
several observables below are built to detect.

\subsection{Lanczos tridiagonalization}
\label{sec:lanczos}

Given a unit-norm starting vector \(q_1\), \(m\) steps of Lanczos build an
orthonormal basis \(Q_m = [q_1 \mid \cdots \mid q_m]\) of the Krylov
subspace \(\mathcal{K}_m(H, q_1) = \mathrm{span}\{q_1, Hq_1, \dots,
H^{m-1}q_1\}\) in which \(H\) restricted to that subspace is tridiagonal:
\[
Q_m^\top H Q_m \;=\; T_m \;=\;
\begin{pmatrix}
\alpha_1 & \beta_1 & & \\
\beta_1 & \alpha_2 & \beta_2 & \\
 & \ddots & \ddots & \ddots \\
 & & \beta_{m-1} & \alpha_m
\end{pmatrix}.
\]
The basis and the tridiagonal entries are produced together by the
three-term recurrence
\[
\alpha_i = q_i^\top H q_i, \qquad
w_i = Hq_i - \alpha_i q_i - \beta_{i-1} q_{i-1}, \qquad
\beta_i = \lVert w_i \rVert, \qquad
q_{i+1} = w_i / \beta_i,
\]
with \(\beta_0 q_0 \equiv 0\). Each step costs exactly one HVP. In finite
precision, orthogonality among the \(q_i\) decays after a few dozen steps,
so \(m\) in the 50--100 range requires \emph{full reorthogonalization} ---
each new \(w_i\) is Gram--Schmidt'd against every previously computed
\(q_1,\dots,q_i\), not just the last one or two -- to keep \(T_m\) an
accurate reduction of \(H\).

\subsection{Stochastic Lanczos Quadrature}
\label{sec:slq}

Diagonalize the small \(m\times m\) matrix \(T_m = U \Theta U^\top\)
(cheap: \(O(m^2)\)--\(O(m^3)\), on the CPU). Gauss quadrature theory says
that for any sufficiently smooth scalar function \(f\),
\[
q_1^\top f(H) q_1 \;\approx\; \sum_{j=1}^m \tau_j^2\, f(\theta_j),
\qquad \tau_j := U_{1j},
\]
where the \(\theta_j\) (eigenvalues of \(T_m\), the \emph{Ritz values}) are
the quadrature \emph{nodes} and \(\tau_j^2\) (squared first components of
\(T_m\)'s eigenvectors) are the quadrature \emph{weights}; \(\sum_j
\tau_j^2 = 1\) since \(U\) is orthogonal. This is exact whenever \(H\)
restricted to \(\mathcal{K}_m(H,q_1)\) captures the relevant spectral
content (e.g.\ exactly, for any \(m \ge 1\), when \(f\) is affine -- see
the trace estimator in Section~\ref{sec:trace}); more generally it is the
degree-\((2m-1)\) moment-matching approximation to \(\int f\,d\mu_{q_1}\)
for the measure \(\mu_{q_1}\) induced by \(q_1\).

A single probe only sees \(H\) through one starting direction. Stochastic
Lanczos Quadrature averages over \(k\) independent random probes
\(q_1^{(1)}, \dots, q_1^{(k)}\), each a normalized Rademacher vector (i.i.d.\
\(\pm1\) entries). Since \(\mathbb{E}[vv^\top] = I\) for such \(v\),
\[
\frac{1}{k}\sum_{l=1}^k \bigl(q_1^{(l)}\bigr)^{\!\top} f(H)\, q_1^{(l)}
\;\xrightarrow[k\to\infty]{}\;
\frac{1}{n}\operatorname{tr} f(H)
\;=\; \int f(\lambda)\, d\mu_H(\lambda).
\]
Each probe \(l\) contributes a discrete measure \(\hat\mu^{(l)} =
\sum_{j} \bigl(\tau_j^{(l)}\bigr)^2 \delta\bigl(\lambda -
\theta_j^{(l)}\bigr)\); pooling all \(k\) probes with weight \(1/k\) each
gives one atomic measure
\(\hat\mu = \frac1k\sum_{l=1}^k \hat\mu^{(l)}\)
that approximates \(\mu_H\). Every observable in Section~\ref{sec:observables}
is a functional of this pooled \((\theta, w)\) atom set.

\subsection{Kernel density estimate for the plotted spectrum}

For the plotted density curve (not the scalar observables -- see
Section~\ref{sec:pooling}), each atom is replaced by a Gaussian kernel of
bandwidth \(\sigma\):
\[
\hat\phi(t) \;=\; \frac{1}{k}\sum_{l=1}^k \sum_{j=1}^m
\bigl(\tau_j^{(l)}\bigr)^2\,
\mathcal{N}\!\bigl(t;\, \theta_j^{(l)},\, \sigma^2\bigr),
\qquad
\mathcal{N}(t;\mu,\sigma^2) = \frac{1}{\sigma\sqrt{2\pi}}
\exp\!\left(-\frac{(t-\mu)^2}{2\sigma^2}\right).
\]

% ============================================================
\section{Implementation: how \code{eigenthings.py} applies the formalism}
\label{sec:implementation}

\subsection{The \code{Eigenthings} class}

\code{Eigenthings.__init__} takes the model, a fixed \code{loss} tensor
(already the result of one forward pass on a batch), and the parameter list
\code{params}. It computes the first-order gradient once,
\begin{quote}
\code{self.grads = torch.autograd.grad(loss, params, create_graph=True, allow_unused=True)}
\end{quote}
This is the \(g(\theta)\) of Section~\ref{sec:hvp}, kept as a differentiable
graph node. Because \(\mathcal{L}\) and hence \(g\) are fixed for the
duration of one spectral snapshot, computing \(g\) once in
\code{__init__} and reusing it across every \code{hvp()} call (rather than
re-walking the whole backward pass per HVP) is what makes the many HVPs
a Lanczos run needs affordable.

\code{hvp(q)} implements Section~\ref{sec:hvp} on a flattened probe vector
\(q \in \R^n\): \code{unpack_vec} reshapes \(q\) into one tensor per
parameter (matching each \(p\)'s shape), \(g^\top v = \sum_p (g_p \cdot
v_p).\mathrm{sum}()\) is accumulated, and a second \code{torch.autograd.grad}
call differentiates that scalar back through \code{self.params} to produce
\(Hv\) (again per-parameter, then re-flattened by \code{pack_list}). Missing
gradients (\code{allow_unused=True}, for parameters the loss doesn't
depend on) are treated as zero.

\code{lanczos(v0)} runs the recurrence of Section~\ref{sec:lanczos} for
\code{self.m} steps: \code{q}/\code{q_prev} track \(q_i, q_{i-1}\);
\code{alpha}/\code{beta} accumulate \(\alpha_i,\beta_i\); \code{Q} keeps
every basis vector generated so far. When \code{self.orthogonalize} is set
(the production default), each new \(w\) is explicitly reduced against
\emph{every} vector in \code{Q}, not just the most recent one --- the full
reorthogonalization called out in Section~\ref{sec:lanczos}. The loop
terminates early if \(\beta_i < 10^{-12}\) (the Krylov subspace has become
invariant --- no new direction to explore), so the realized depth \(k =
\code{len(alpha)}\) can be less than \code{self.m}. The tridiagonal \(T_m\)
is assembled explicitly from \code{alpha}/\code{beta} and returned.

\code{lanczos_quadrature(v0)} calls \code{lanczos}, moves the (small) \(T_m\)
to CPU for \code{torch.linalg.eigh}, and reads off nodes
\(\theta_j = \code{eigvals}\) and weights \(\tau_j^2 =
\code{eigvecs[0, :]**2}\) exactly as in Section~\ref{sec:slq}.

\code{density_from_lanczos(nodes, weights, t_grid, sigma)} evaluates the
KDE sum of Section~\ref{sec:formalism} on an arbitrary grid, skipping atoms
whose weight is negligible (\(|\tau_j^2| \le 10^{-14}\)) for efficiency.

\subsection{Probe sampling: \code{rademacher_probe}}

\code{rademacher_probe} draws \(\pm1\) entries via a CPU
\code{torch.Generator} and only moves the result to the target device
afterwards, because CUDA/MPS generators are not bitwise-interchangeable
with the CPU one --- this keeps the probe sequence identical regardless of
accelerator, which matters for reproducibility across runs/hardware.

\subsection{\code{estimate_density}: orchestrating one spectral snapshot}
\label{sec:estimate_density}

\code{estimate_density(model, m, k, sigma_frac, loss, probe_seed, \ldots)}
is the entry point used once per logged training step. It runs in two
phases:

\begin{enumerate}
\item \textbf{Cheap top-eigenvalue pre-pass.} A single Rademacher probe
(seeded from \code{probe_seed}, its own generator so it does not disturb
the sequence used below) is run through \code{lanczos_quadrature} with a
shallow depth \code{top_eig_probe_m} (default 20 -- see
Section~\ref{sec:accuracy}, \code{top_eig} converges by \(m\ge20\)). The
largest node gives a cheap estimate of \(|\lambda_{\max}|\). This exists
solely to set the absolute KDE bandwidth: \(\sigma =
\max(\code{sigma_frac} \cdot |\lambda_{\max}|,\, \code{eps})\), a
\emph{fraction} of \(\lambda_{\max}\) rather than a fixed absolute value,
since \(\lambda_{\max}\) itself can move by orders of magnitude over a
training run and a fixed \(\sigma\) would make density plots incomparable
across checkpoints.
\item \textbf{Main \(k\)-probe, \(m\)-step pass.} The same
\code{Eigenthings} instance (reusing its once-computed \code{grads}) is
reconfigured to depth \code{m}, and a fresh generator reseeded from the
\emph{same} \code{probe_seed} draws \(k\) Rademacher probes. Calling
\code{estimate_density} repeatedly with the same \code{probe_seed}
(e.g.\ once per checkpoint of a training run) therefore draws bit-identical
probes every time, isolating parameter-driven spectral change across
checkpoints from probe-sampling noise.
\end{enumerate}
The plotting grid \code{t_grid} spans \([\min(\text{nodes}) - 6\sigma,\;
\max(\text{nodes}) + 6\sigma]\) at step \(\sigma/5\), clamped to
\([200, 50000]\) points so that a very wide node range (e.g.\ right after
the top-eig pre-pass makes \(\sigma\) very small) cannot silently blow the
grid past a usable size. \code{SpectralEstimate} bundles both outputs: the
KDE-smoothed \code{density}/\code{t_grid} (plotting only) and the raw,
un-smoothed per-probe \code{probe_nodes}/\code{probe_weights}
(everything in Section~\ref{sec:observables}).

\subsection{Integration into the training loop}

\code{training/loop.py} calls \code{estimate_density} on a schedule
(\code{SpectralSchedule}, geometric-then-dense around the
train/test-accuracy-gap transition), always against the \emph{same} fixed
batch (\code{spectral_batch_size}, default 512, sampled once before
training starts) so that, per snapshot, only \(\theta\) differs and neither
the data nor the probes are confounds. The raw nodes/weights returned by
\code{estimate_density} are passed straight into
\code{compute_spectral_observables_with_stderr}
(Section~\ref{sec:observables}), and both the point estimates and their
stderrs are logged into the training history per checkpoint.

% ============================================================
\section{Observables}
\label{sec:observables}

All eight observables below are implemented in
\code{analysis/spectral_observables.py}, \code{compute_spectral_observables},
and are computed from the \emph{raw} pooled Ritz atoms, not the
KDE-smoothed density --- the smoothing bandwidth would otherwise bias
tail-sensitive quantities such as \code{top_eig} and outlier detection.

\subsection{Pooling probes into one discrete measure}
\label{sec:pooling}

\code{_pool_probes} concatenates the \(k\) probes' \((\theta^{(l)},
\tau^{(l)2})\) pairs into one atom set, dividing each probe's weights by
\(k\) so the pooled weights sum to 1 overall (each probe contributes total
mass \(1/k\), matching the SLQ average of Section~\ref{sec:slq}):
\[
\{(\theta_j, w_j)\}_{j=1}^{N}, \qquad N = \sum_{l=1}^k (\text{realized depth of probe } l),
\qquad w_j = \tau_j^{(l)2}/k.
\]
Below, \(\lambda_{\max} := \theta_N\) after sorting ascending (the largest
pooled node), and \(n\) denotes \code{n_params}, the true parameter count
--- \emph{not} \(N\), which scales with \(m\) and \(k\) and would make
several observables below drift with Lanczos depth/probe count rather than
track the actual spectrum.

\subsection{\code{top_eig}}
Mathematically, the estimate of \(\lambda_{\max}(H)\). Code: the largest
pooled node, \code{nodes[-1]} after sorting. Precise because it is read off
the raw atoms rather than the KDE curve, which would blur/underestimate the
tail.

\subsection{\code{bulk_edge}}
A robust threshold separating the low-curvature ``bulk'' of the spectrum
from outlier eigenvalues, defined as a weighted median plus a multiple of
the weighted median absolute deviation (MAD):
\[
\code{bulk_edge} \;=\; \mathrm{median}_w(\theta) + c\cdot 1.4826\cdot
\mathrm{MAD}_w(\theta), \qquad c = \code{bulk_mad_multiplier} = 5,
\]
where \(\mathrm{median}_w\) and \(\mathrm{MAD}_w = \mathrm{median}_w(|\theta -
\mathrm{median}_w(\theta)|)\) are computed against the \emph{weighted} CDF
of the pooled atoms (\code{_weighted_quantile}, linear interpolation of
\(q\) against \(\mathrm{cdf}_j = (\sum_{i\le j} w_i - \tfrac12 w_j)/\sum_i
w_i\), i.e. each atom's cumulative-probability coordinate sits at its own
mass's midpoint), not a plain \code{numpy.median} over the raw node values.
The \(1.4826\) factor rescales MAD to be consistent with a Gaussian
standard deviation. This weighting is not cosmetic: Section~\ref{sec:accuracy}
documents that the unweighted version was found to be biased 30--50\% on a
spectrum with a realistic bulk/outlier imbalance, because Lanczos node
placement is moment-matching, not density-matching -- nodes cluster near
spectral edges regardless of how much true probability mass sits there, so
treating every node as equally likely (unweighted) systematically
overweights the edges.

\subsection{\code{outlier_count}}
The number of \emph{distinct} eigenvalues sitting above \code{bulk_edge},
as opposed to the number of Ritz values above it (Lanczos rediscovers the
same true outlier eigenvalue as several near-identical Ritz values across
different probes/depths). Code: take pooled nodes with \(\theta_j >
\code{bulk_edge}\), sort, and count runs separated by gaps larger than
\(\code{resolution} = \max(\code{resolution_frac}\cdot|\lambda_{\max}|,
\code{eps})\) (default \(\code{resolution_frac}=0.05\)) -- i.e.\ the number
of gaps exceeding \(\code{resolution}\), plus one.

\subsection{\code{trace}}
\label{sec:trace}
The Hutchinson-via-Lanczos trace estimator,
\[
\widehat{\operatorname{tr}}(H) \;=\; n\sum_{j=1}^N w_j\,\theta_j.
\]
This follows directly from Section~\ref{sec:slq} with \(f(\lambda) =
\lambda\): \(\mathbb{E}_v[v^\top H v] = \operatorname{tr}(H)\) for Rademacher
\(v\), and \(q_1^\top H q_1 = \sum_j \tau_j^2 \theta_j\) is exact for
\emph{any} \(m \ge 1\) --- it is the first-moment identity of Gauss
quadrature (in fact \(q_1^\top H q_1 = \alpha_1\), the very first Lanczos
coefficient). This is why, uniquely among the eight fields,
\code{trace}'s accuracy is governed entirely by \(k\) (Monte Carlo/Hutchinson
variance, \(O(k^{-1/2})\)) and is essentially insensitive to \(m\).

\subsection{\code{spectral_entropy} and \code{effective_rank}}
The Shannon entropy of a \emph{magnitude-reweighted} spectral measure. Model
each of the \(n\) exact eigenvalues \(\lambda_i\) as carrying uniform prior
mass \(1/n\), reweighted by \(|\lambda_i|\) (so directions with negligible
curvature, positive or negative, don't count toward ``effective rank''):
\[
p_i \;=\; \frac{|\lambda_i|}{n Z}, \qquad
Z := \frac1n\sum_{i=1}^n |\lambda_i| \;\approx\; \sum_{j=1}^N w_j |\theta_j|
\quad\text{(pooled SLQ estimate of } \textstyle\frac1n\sum_i|\lambda_i|\text{)}.
\]
Then
\[
S \;=\; -\sum_{i=1}^n p_i \log p_i
\;=\; \log(nZ) \;-\; \frac{1}{nZ}\sum_{i=1}^n |\lambda_i|\log|\lambda_i|
\;\approx\; \log(nZ) \;-\; \frac{1}{Z}\sum_{j=1}^N w_j |\theta_j| \log|\theta_j|,
\]
using the same pooled-SLQ substitution for \(\frac1n\sum_i(\cdot)\) in the
second term. This is exactly \code{spectral_entropy} in the code
(\code{z = sum(weighted_abs)}, \code{cross_term = sum(weighted_abs *
log_abs)}, \code{spectral_entropy = log(n_params*z) - cross_term/z}),
computed in closed form off the raw atoms with no histogram binning, so
there is no artificial cap on the result. \code{effective_rank} \(=
e^S\) is the ``perplexity'' of \(p\): the number of eigen-directions
(out of \(n\)) carrying non-negligible \(|\)curvature\(|\), regardless of
sign. Because \(n\) (the true parameter count) rather than \(N\) (the
pooled Ritz-node count, which scales with \(mk\)) appears in the \(\log(nZ)\)
term, \(S\) and \(e^S\) are comparable across different \(m,k\) settings
rather than drifting with them.

\subsection{\code{negative_mass}}
The probability mass at eigenvalues meaningfully below zero,
\emph{margin-thresholded} rather than a raw \(\lambda<0\) count:
\[
\code{negative_mass} \;=\; N_\epsilon \;=\; \sum_{j:\, \theta_j < -\epsilon} w_j,
\qquad \epsilon = \code{margin_c}\cdot|\lambda_{\max}|,\ \ \code{margin_c}=0.05.
\]
The margin exists because Lanczos produces spurious small-magnitude
negative Ritz values immediately around a true eigenvalue near zero; a raw
\(<0\) threshold counts that numerical noise as curvature evidence (see
Section~\ref{sec:accuracy} for the measured effect). \code{margin_c} shares
its default with \code{resolution_frac} (both are ``a small fraction of
\(|\lambda_{\max}|\)'' length scales) but is an independently tunable knob:
one sets a spectral resolution for clustering, the other a decision
boundary for ``genuinely negative.''

\subsection{\code{conditioning}}
A cheap proxy for how far the extreme eigenvalue sits above the bulk,
\[
\code{conditioning} = \frac{\code{top_eig}}{\code{bulk_edge}},
\]
undefined (\code{NaN}) if \(|\code{bulk_edge}| \le \code{eps}\). Chosen over
a fixed absolute threshold or a Marchenko--Pastur-style edge (which needs
an aspect ratio not well-defined for a Hessian).

\subsection{Error bars: \code{stderr} over probes}
\code{compute_spectral_observables_with_stderr} returns
\((\text{point}, \text{stderr})\): \code{point} pools all \(k\) probes (as
above, the least-biased estimate); \code{stderr} recomputes every
observable \(k\) times, once per single-probe atom set, and reports the
standard error of the mean across those \(k\) values,
\(\mathrm{sd}(\{x^{(l)}\}_{l=1}^k)/\sqrt{k}\) with \(\mathrm{ddof}=1\). This
costs no extra HVPs (it reuses the already-computed per-probe atoms) and is
a pragmatic uncertainty proxy rather than the literal standard error of
every point estimate -- e.g.\ \code{top_eig} is a pooled \emph{max}, not a
mean of per-probe maxes, so the point estimate is not the mean of the
values whose spread \code{stderr} reports.

% ============================================================
\section{Accuracy: how \texorpdfstring{\(m\)}{m} and \texorpdfstring{\(k\)}{k} govern each observable}
\label{sec:accuracy}

\subsection{What \(m\) and \(k\) each control}

\(m\) (Lanczos steps per probe, \code{spectral_m}) sets the depth of the
Krylov subspace built per probe, i.e.\ \emph{algorithmic} convergence: how
well the small \(T_m\)'s Ritz values approximate the true eigenvalues
\(H\) has along that probe's Krylov direction. \(k\) (probe count,
\code{spectral_k}) sets how many independent Rademacher directions are
averaged, i.e.\ \emph{Monte Carlo} convergence of the Hutchinson-style
average in Section~\ref{sec:slq}, governed by the CLT (\(O(k^{-1/2})\)
error). These are different axes, and --- this is the central empirical
finding of \code{reports/spectral_validation.md} --- different observables
are bottlenecked by different axes, so no single \((m,k)\) is simultaneously
``cheapest'' for every field. Total cost per spectral snapshot is
\(\approx m \cdot k\) HVPs (one HVP per Lanczos step, per probe), and this
snapshot recurs at every logged step of a training run, so \((m,k)\) is the
dominant lever on the recurring cost of \code{run_spectral=True}.

\subsection{Per-observable convergence}

The table below summarizes the sweeps in \code{reports/spectral_validation.md}
against two exact-reference fixtures: a $\sim$1026-parameter random-init
MLP with a materialized Hessian (\code{tests/reference_hessian.py}), and an
engineered quadratic-form fixture with the same parameter count but a
production-like spectral dynamic range (exact \code{top_eig}\(=150\),
\code{bulk_edge}\(=1.109\), \code{conditioning}\(=135.3\), i.e.\ \(\sim\)31\(\times\)
wider than the first fixture), built specifically to exercise
\code{bulk_edge}/\code{conditioning}, which the near-flat first fixture
could not. \emph{The qualitative governing axis (\(m\) vs.\ \(k\)) is a
property of the SLQ algorithm and should generalize; the specific numeric
thresholds are shaped by each fixture's spectral gap and should not be
copy-pasted onto a production run without their own check.}

\begin{center}
\begin{tabular}{@{}l l l l@{}}
\toprule
Observable & Governing axis & Min. \(m\) & Min. \(k\) \\
\midrule
\code{trace} & \(k\) (\(m\) irrelevant) & 1 & \(\ge100\) (\(\ge300\) tighter) \\
\code{top_eig} & \(m\) & \(\ge20\) & 1 \\
\code{negative_mass} (margin-thresh.) & \(m\) & \(\ge30\) (50 for full margin) & \(\ge30\) \\
\code{spectral_entropy} & \(m\) & \(\ge20\) (30 w/ margin) & \(\ge30\) \\
\code{effective_rank} & \(m\) & \(\ge20\) (30 w/ margin) & \(\ge30\) \\
density (plotted curve) & \(m\), then \(k\) & \(\ge20\) (30 w/ margin) & \(\ge100\) \\
\code{bulk_edge} (post-fix) & \(m\) & \(\ge75\) (100 for \(<\)0.2\%) & not swept \\
\code{conditioning} (post-fix) & \(m\) & \(\ge75\) (100 for \(<\)0.2\%) & not swept \\
\bottomrule
\end{tabular}
\end{center}

Notable points behind the table:
\begin{itemize}
\item \textbf{\code{trace}} converges purely in \(k\) because, as shown in
Section~\ref{sec:trace}, the underlying quadrature identity is exact in
\(m\) for the first moment; all error is Hutchinson/Monte-Carlo variance.
Measured RMS relative error (\(R=20\) independent repeats) was 21.0\% at
\(k{=}10\), 4.6\% at \(k{=}100\), 2.2\% at \(k{=}1000\) --- \(\sqrt{k}\)
convergence confirmed at a fitted log-log slope of \(-0.511\) against a
\(-0.5\) theoretical rate. Reaching \(\sim\)1\% relative error would need
\(k\) in the low thousands.
\item \textbf{\code{negative_mass}} was originally defined with a raw
\(\theta<0\) threshold and needed \(m \ge 200\) to converge (non-monotonically
in \(m\)) --- traced to 193 of 1026 exact eigenvalues on the reference
fixture sitting within 5\% of \(|\lambda_{\max}|\) of zero, a band where
Lanczos needs substantial depth just to resolve the \emph{sign} correctly.
Switching to the margin-thresholded definition of
Section~\ref{sec:observables} (excluding that band by construction) dropped
the requirement an order of magnitude, to \(m\ge30\)--50.
\item \textbf{\code{bulk_edge}/\code{conditioning}} were the deepest
\(m\)-governed fields once correctly estimated, but their original
implementation (plain, \emph{unweighted} \code{numpy.median}/MAD over the
pooled Ritz nodes) was found to be 30--50\% biased \emph{at every \(m\)
tested, including \(m{=}300\)}, with negligible probe-to-probe spread --
the signature of a biased estimator, not under-convergence, since raising
\(m\) further would not have fixed it. This was an estimator bug (Section
\ref{sec:observables}'s weighted-median fix), not a parameter-tuning
problem; only after the fix does the \(m\ge75\)--100 threshold in the table
reflect genuine Lanczos-depth convergence.
\item \textbf{\code{spectral_entropy}/\code{effective_rank}} converge an
order of magnitude faster in \(m\) than \code{negative_mass} despite
reading the same pooled \(|\lambda|\) distribution, because the entropy's
closed form (Section~\ref{sec:observables}) only needs Ritz mass to land in
roughly the right coarse bin, whereas \code{negative_mass} needs the exact
sign right at a zero crossing --- a much finer distinction that coarse
convergence can't average over.
\end{itemize}

\subsection{Cost and production defaults}

Current production defaults, \code{configs/base.py}:
\[
\code{spectral_m} = 50, \qquad \code{spectral_k} = 100,
\]
i.e.\ \(\approx 5{,}000\) HVPs per logged spectral snapshot. Cross-referencing
against the table above: \(k{=}100\) is a deliberate compromise for
\code{trace} (\(\sim\)5\% relative error, not the \(\sim\)1\% original
target, since that would need \(k\) in the low thousands and \(k\) multiplies
snapshot cost directly for a quantity logged repeatedly through training,
not once); \(m{=}50\) comfortably clears the \(m\ge20\)--30 threshold shared
by \code{top_eig}, \code{negative_mass}, \code{spectral_entropy}, and
\code{effective_rank}, but sits \emph{below} the \(m\ge75\)--100 threshold
the validation doc establishes for \code{bulk_edge}/\code{conditioning} to
reach sub-1.5\%/sub-0.2\% relative error post-fix. (The inline comment on
\code{spectral_m} in \code{configs/base.py} cites ``\(m{=}50\), \(k{=}300\)''
as clearing every \(m\)-bound field, but the checked-in default sets
\code{spectral_k=100}, not 300, and \code{spectral_validation.md}'s own
summary table puts \(m\ge75\) --- not 50 --- as the requirement for
\code{bulk_edge}/\code{conditioning}; this is a live discrepancy between
the code comment and both the configured value and the validation doc's
recommendation, worth resolving directly rather than assuming either side
is authoritative.) \code{bulk_edge}/\code{conditioning} have also not been
swept over \(k\) at all, so their behavior below the $\sim$0.2\% level with
non-default \(k\) is unverified.

Because cost scales as \(m\cdot k\), and \(k\) is set by \code{trace} while
\(m\) is set (once \code{bulk_edge}/\code{conditioning} are the binding
constraint) independently of \(k\), the two knobs should generally be tuned
separately per analysis rather than raised together: if only \code{top_eig}
or the entropy family is needed, \(m\approx20\)--30 with \(k\) as low as 1 is
sufficient and far cheaper than the shared production default; if
\code{trace} accuracy tighter than \(\sim\)5\% is needed for a specific run,
that is a reason to raise \code{spectral_k} for that run rather than the
global default, given the direct cost multiplier on every logged checkpoint.

\end{document}
